\section{Постановка задачи}

В рамках курсовой работы необходимо разработать базу данных для предметной области <<Учёт компьютерных шрифтов в электронной типографии>>.

База данных должна обеспечивать хранение и обработку сведений о шрифтовых ресурсах типографии, рабочих станциях и лицензионных ограничениях. В системе необходимо учитывать гарнитуры, входящие в них начертания, категории шрифтов, технические форматы файлов, поддерживаемые языки и символы, издательства, лицензии, операционные системы, рабочие станции и факты установки начертаний на конкретные станции.

Необходимо выполнить следующие этапы работы:

\begin{itemize}
    \item провести анализ предметной области и определить основные цели функционирования базы данных;
    \item выделить сущности, их атрибуты, связи и роли пользователей системы;
    \item построить ER-диаграмму и схему объектов базы данных;
    \item выполнить логическое проектирование таблиц, выбрать типы данных и ограничения целостности;
    \item реализовать SQL-скрипт создания таблиц базы данных;
    \item подготовить программу заполнения базы данных тестовыми данными;
    \item составить SQL-запросы, демонстрирующие основные пользовательские операции;
    \item получить и проанализировать планы выполнения запросов.
\end{itemize}

Разрабатываемая база данных должна позволять:

\begin{itemize}
    \item хранить каталог гарнитур и начертаний с их типографическими и техническими характеристиками;
    \item фиксировать принадлежность гарнитур к категориям и издательствам;
    \item учитывать форматы файлов начертаний и поддерживаемые языки;
    \item хранить сведения о символах, доступных в конкретных начертаниях;
    \item вести учёт лицензий, сроков их действия, стоимости и ограничения по числу рабочих станций;
    \item хранить сведения о рабочих станциях и установленных на них начертаниях;
    \item выполнять аналитические запросы по форматам, категориям, языкам, символам и установленным шрифтам.
\end{itemize}

Результатом работы должны стать спроектированная и реализованная база данных, заполненная тестовыми данными, набор SQL-запросов для проверки её работы, а также отчёт с описанием аналитики, проектирования, реализации и результатов выполнения запросов.
